\par\bigskip\Needspace{9\baselineskip}
\SourceLeaf{222}
\renewcommand{\BookRunningTitle}{Nachlass. Disquisitiones generales de congruentiis}
\pdfbookmark[0]{Disquisitiones generales de congruentiis}{disquisitiones-generales-de-congruentiis}
\begin{center}
{\large DISQUISITIONES GENERALES DE CONGRUENTIIS.}\par\medskip
ANALYSIS RESIDUORUM CAPUT OCTAVUM.
\end{center}
\BookRule{.18\linewidth}
\renewcommand{\Article}[1]{\par\medskip\Needspace{4\baselineskip}\pdfbookmark[1]{Art. #1}{disquisitiones-generales-art-#1}\centerline{\textbf{#1.}}\nobreak\smallskip}
\Article{330}
Quae in Sectionibus praecedentibus de congruentiis sunt tradita, simplicissimos tantum casus attinent methodisque particularibus plerumque sunt eruta. In hac Sectione periculum faciemus congruentiarum theoriam, quantum quidem adhuc licet, ad altiora principia reducere, simili fere modo ut \emph{aequationum} theoria considerari solet, quacum insignis intercedit analogia, uti iam saepius observavimus. Quoniam igitur omnes congruentiae algebraicae unicam incognitam involventes ad hanc formam reduci possunt
\[
X\equiv0
\]
ubi $X$ est functio algebraica incognitae $x$, nullas fractiones involvens, huiusmodi functiones imprimis erunt considerandae.

\Article{331}
Si $P,Q$ sint functiones indeterminatae $x$ huius formae
\[
\begin{gathered}
A+Bx+Cxx+Dx^3+\ldots\\
H+Ix+Kxx+Lx^3+\ldots
\end{gathered}
\]
(quales abhinc semper per \emph{functiones} simpliciter designamus) et in utraque coëfficientes similium ipsius $x$ potestatum secundum quemcunque modulum sint 
\SourceLeaf{223}
congrui, \emph{functiones secundum hunc modulum congruae} dicentur. Perspicuum autem est, functiones congruas, si pro indeterminata valores aequales aut congrui accipiantur, valores congruos nancisci. Quae in Capp \textsc{i}. et \textsc{ii}. de \emph{numeris} demonstravimus, plerumque etiam de functionibus sunt tenenda; ita si $P\equiv P'$, $Q\equiv Q'$, $R\equiv R'$ etc., patet, fore $P+Q+R\text{ etc.}\equiv P'+Q'+R'+\text{etc.}$; $P-Q\equiv P'-Q'$; $PQ\equiv P'Q'$; $PQR\text{ etc.}\equiv P'Q'R'\text{ etc.}$ Demonstrationes facillimae, possuntque simili modo adornari ut Cap.\ $\mathrm{I}^{\mathrm{mo}}$.

Si $PQ\equiv R$, functionem $Q$ per $\frac RP$ designabimus apposito modulo, dicemusque, $Q$ esse quotientem, si $R$ per $P$ secundum hunc modulum dividatur. Manifestum autem est, loco ipsius $Q$ omnes functiones ipsi congruas accipi posse, quas omnes tamquam \emph{unicum} valorem spectabimus. Infra vero ostendemus, quibus casibus talis quotiens plures valores (i.\,e. incongruos) nancisci possit.

\Article{332}
Si modulus sit numerus primus et divisor $Q$ unicum tantum terminum involvat $Hx^h$, cuius coëfficiens $H$ per modulum non dividitur, i.\,e. si modo $H$ non sit $\equiv0$, quotiens plures valores habere nequit. Si enim esset $QA\equiv P$ et $QB\equiv P$, foret $Q(A-B)\equiv0$. Iam sit
\[
Q\equiv\dotstwo+Hx^h+Ix^{h+1}+\text{etc.}
\]
ita ut $H$ per $p$ non dividatur, et
\[
A-B\equiv Lx^l+Mx^{l+1}+\text{etc.}
\]
ita ut $L$ per $p$ non dividatur (hanc autem formam $A-B$ habebit, quia supponimus $A$ non $\equiv B$). Foretque $Q(A-B)\equiv HLx^{h+l}+\text{etc.}\equiv0$. Q.\,E.\,A., quia $HL$ non $\equiv0$.

Facile iam regulae dantur functionem $P$ per $Q$, siquidem fieri potest, dividendi; sit
\[
\begin{aligned}
P&\equiv ax^\alpha+bx^{\alpha+1}+cx^{\alpha+2}+\text{etc.}+kx^\kappa\\
Q&\equiv mx^\mu+nx^{\mu+1}+qx^{\mu+2}+\text{etc.}+tx^\tau
\end{aligned}
\]
ita ut $a,k,m,t$ per modulum non dividantur, debetque esse $\alpha$ non $<\mu$, $\kappa$ non $<\tau$. Divisio autem simili modo institui potest, ut in calculo logistico communi, modo semper pro quotiente numerus integer accipiatur; scilicet quotiens semper
\SourceLeaf{224}
hanc formam habebit $\frac rm$, quod secundum modulum determinari debet. Iam si postquam $\kappa+\mu-\alpha-\tau+1$ termini sunt inventi, residuum remaneat, quod erit formae
\[
Ax^{\kappa+\mu-\tau+1}+Bx^{\kappa+\mu-\tau+2}+\dotstwo+Cx^\kappa
\]
neque omnes coëfficientes $A,B,C\dotstwo$ sint $\equiv0$, $P$ per $Q$ dividi nequit.

Ceterum patet, divisionem etiam a terminis, qui maximas dimensiones habent, $kx^\kappa$, $tx^\tau$ incipi potuisse; operatio facilitabitur, si $Q$ ad formam redigatur
\[
mx^\mu(1+qx+rxx+\text{etc.})
\]
unde fiet posito $mv\equiv1$
\[
\frac PQ\equiv\frac{vP:x^\mu}{1+qx+\text{etc.}}
\]
tunc vero divisio per methodos communes perfici potest.

\Article{333}
\textsc{Theorema.} \emph{Si $x\equiv a$ fuerit radix congruentiae $\xi\equiv0$, $\xi$ per $x-a$ dividi poterit secundum congruentiae modulum.}

\emph{Demonstratio.} Si enim dividi non posset, foret $\xi\equiv(x-a)\xi'+b$, ita ut $b$ per modulum dividi non posset. Iam si $x$ ponatur $\equiv a$, $\xi$ fiet $\equiv0$ (hyp.), quare $(x-a)\xi'+b\equiv0$; sed tunc etiam $(x-a)\xi'\equiv0$, quare $b$ necessario erit $\equiv0$.

\Article{334}
\textsc{Problema.} \emph{Datis binis functionibus, earum communem divisorem (maximae dimensionis) invenire secundum modulum datum.}

\emph{Solutio.} Sint functiones $A,B$. Habeat $A$ totidem aut plures dimensiones quam $B$; dividatur $A$ per $B$, si fieri potest sine residuo, $B$ erit divisor communis quaesitus. Si residuum maneat $C$, hoc inferiorem dimensionem habebit, quam $B$. Sit itaque
\[
A\equiv aB+C,\quad B\equiv bC+D,\quad C\equiv cD+E,\quad\text{etc.}
\]
ita ut $A,B,C,D,a,b,c$ etc. sint functiones, et dimensiones functionum $A,B,C,D$ etc. constituant seriem decrescentem. Iam si tandem aliqua divisio succedat, ex.\ gr.\ $D\equiv dE$, ultimus divisor erit divisor communis quaesitus; si vero nulla succedat, tandem ad residuum pervenietur, quod nullam dimensionem
\SourceLeaf{225}
habeat i.\,e. ad \emph{numerum}; hoc autem casu functiones $A,B$ communem divisorem non habent.

\emph{Demonstr.} Si divisor $E$ functionem praecedentem sine residuo dividat, omnes antecedentes dividere facile perspicitur; quare $E$ erit divisor communis functionum $A,B$. Q.\,E.\,Pr. Si autem daretur divisor maioris dimensionis, puta $E'$, hic propter $C\equiv A-aB$ etiam $C$ similique argumento etiam $D$ etc. adeoque $E$ divideret, functio maioris dimensionis functionem minoris. Q.\,E.\,A.\ Q.\,E.\,Scd. Hinc etiam patet, si divisor communis ullius dimensionis datur, ad residuum nullius dimensionis perveniri non posse; alias enim functio nullius dimensionis per functionem alicuius dimensionis divideretur. Q.\,E.\,A.

\Article{335}
\textsc{Theorema.} \emph{Si $A,B$ sint functiones inter se primae secundum modulum $p$; $A$ autem dimensionis $\alpha$, $B$ dimensionis $\beta$; inveniri poterunt functiones $P,Q$, dimensionum quae sunt respective $<\beta$, $<\alpha$, ita ut}
\[
PA+QB\equiv1\;(\mathit{mod.}\,p)
\]

\emph{Demonstr.} Hoc enim casu erit
\[
A\equiv aB+C,\quad B\equiv bC+D,\quad\text{etc.}\quad K\equiv kL+M
\]
ita ut dimensiones functionum $A,B,C,D,\dotstwo K,L,M$ continuo decrescant et $M$ nullam dimensionem habeat. Iam formentur series
\[
\begin{gathered}
a,\quad a',\quad a'',\quad a''',\ldots a^{(\kappa)}\\
1,\quad b,\quad b',\quad b'',\ldots b^{(\kappa-1)}
\end{gathered}
\]
ita ut
\[
\begin{array}{llll}
a'\equiv ba+1&a''\equiv ca'+a&a'''\equiv da''+a'&\text{etc.}\\
b'\equiv cb+1&b''\equiv db'+b&b'''\equiv eb''+b'&\text{etc.}
\end{array}
\]
eritque
\[
A-aB\equiv+C,\quad bA-a'B\equiv-D,\quad b'A-a''B\equiv+E,\quad\text{etc.}
\]
uti sine negotio perspicitur; hinc tandem
\[
b^{(\kappa-1)}A-a^{(\kappa)}B\equiv\pm M
\]
Iam sit $\frac1{\pm M}\equiv\mu$, eritque ponendo $P\equiv\mu b^{(\kappa-1)}$, $Q\equiv-\mu a^{(\kappa)}$
\SourceLeaf{226}
\[
PA+QB\equiv1
\]
Porro vero manifestum est,
\[
\begin{aligned}
\text{Dimens. ipsius }B+\text{Dim. ipsius }a&\text{ esse }=\text{Dim. }A.\\
\text{Dim. }C+\text{Dim. }b&=\pdot\text{Dim. }B\\
&\text{etc.}\\
\text{Dim. }L+\text{Dim. }k&=\text{Dim. }K
\end{aligned}
\]
Quare
\[
\text{Dim. }L+\text{Sum. Dim. }a,b,\dotstwo k=\text{Dim. }A
\]
Patet vero dimensionem ipsius $a^{(\kappa)}$ adeoque etiam
\[
\text{Dim. ipsius }Q\text{ esse }=\text{Sum. Dim. }a,b,c,\dotstwo\text{ i. e. }=\alpha-\text{Dim. }L
\]
itaque
\[
\text{Dim. ipsius }P=\beta-\text{Dim. }L\qquad\text{Q. E. D.}
\]

\Article{336}
Hinc autem sequitur, si $M$ est divisor communis maximae dimensionis functionum $A,B$, semper poni posse
\[
AP+BQ\equiv M
\]

Exempla praecedentis theorematis brevitatis gratia omitto, sed lectores non negligent, per ea facilitatem huius generis problemata tractandi sibi comparare. Ceterum operae pretium erit admonere, theorema praecedens etiam de functionibus absolute sumtis valere, quarum quidem coëfficientes sint numeri rationales. Hoc ex demonstrationis modo per se elucebit. Nobis autem ei rei immorari non licet. Similia lector etiam non admonitus in sequentibus observabit.

Si $A$ nec cum $B$ nec cum $C$ divisorem ullius dimensionis communem habeat, etiam cum producto $BC$ nullum habebit divisorem communem. Sit enim
\[
PA+QB\equiv1,\quad\text{erit}\quad PAC+QBC\equiv C
\]
Iam si $A$ cum $BC$ divisorem $M$ communem haberet, hic etiam ipsam $C$ divideret contra hyp. Hinc generaliter si functio $A$ ad $B,C,D$ etc. prima, etiam ad omnium productum erit prima.

Si $A,B,C,D$ etc. nullum divisorem habeant omnibus communem, fieri potest
\SourceLeaf{227}
\[
PA+QB+RC+SD+\text{etc.}\equiv1
\]
Sit divisor maximae dimensionis inter $A$ et $B,M$; inter $M$ et $C,M'$; inter $M'$ et $D,M''$ etc.: patet, ultimum huius seriei terminum fore nullius dimensionis (hyp.). Quare poni poterit
\[
aA+bB\equiv M,\quad mM+cC\equiv M',\quad m'M'+dD\equiv M'',\quad\text{etc.}
\]
unde substitutionibus factis theorematis veritas apparet.

\Article{337}
\textsc{Theorema.} \emph{Si $A,B,C$ etc. sint functiones inter se primae (quarum binae quaeque nullum habeant divisorem communem) secundum modulum $p$, et functio $M$ secundum eundem modulum per singulas sit divisibilis; etiam per omnium productum erit divisibilis.}

\emph{Demonstr.} Poni enim potest $PA+QB\equiv1$, quare erit
\[
\frac MA Q+\frac MB P\equiv\frac M{AB}
\]
Iam quum $C$ ad $AB$ prima, erit etiam $M$ per $ABC$ divisibilis similique ratiocinio per $ABCD$ etc.

\Article{338}
Si congruentia $\xi\equiv0$ habeat radices $x\equiv a$, $x\equiv b$, $x\equiv c$ etc., $\xi$ per productum ex $(x-a)$, $(x-b)$, $(x-c)$ etc. dividi poterit; cum enim $a,b,c\pdot$ etc. supponantur incongrui, functiones $x-a,x-b,x-c$ etc. erunt primae inter se, et quum $\xi$ per singulas dividatur, etiam per productum ex omnibus dividetur. Hinc patet, radicum multitudinem congruentiae dimensionem superare non posse: quae est demonstratio huius theorematis, quam polliciti sumus.

Sed simul hinc perspicitur, quomodo congruentiarum solutio partem tantummodo constituat multo altioris disquisitionis, scilicet de resolutione functionum in factores. Manifestum est, congruentiam $\xi\equiv0$ nullas habere radices reales, si $\xi$ nullos factores unius dimensionis habeat; at hinc nihil obstat, quominus $\xi$ in factores duarum, trium pluriumve dimensionum resolvi possit, unde radices quasi \emph{imaginariae} illi attribui possint. Revera, si simili licentia, quam recentiores mathematici usurparunt, uti talesque quantitates imaginarias introducere 
\SourceLeaf{228}
voluissemus, omnes nostras disquisitiones sequentes incomparabiliter contrahere licuisset; sed nihilominus maluimus omnia ex principiis deducere\begingroup\renewcommand{\thefootnote}{*)}\footnote{Alia forsan occasione de hac re opinionem nostram fusius explicabimus.}\endgroup.

\Article{339}
\emph{Functiones} secundum modulum determinatum \emph{primae} vocantur, quae per nullas functiones inferiorum dimensionum secundum hunc modulum dividi possunt.

Ita omnes functiones unius dimensionis erunt primae, functiones autem duarum dimensionum aut erunt primae aut ex binis unius dimensionis compositae: quare $\xi$ erit functio prima duarum dimensionum, si congruentia $\xi\equiv0$ nullas radices reales admittit. Ex. gr. $xx+x+1$ pro modulo $5$ est prima, quia
\[
xx+x+1\equiv(x-2)^2-3\;(\mathrm{mod.}\,5)
\]
et $3$ non-residuum quadraticum numeri $5$.

Hae vero functiones primae prae omnibus attentionem nostram desiderant. Quamvis enim aliae quam primi gradus ad inveniendas radices reales inservire non possint, amplior earum consideratio tum ob insignes ipsarum proprietates tum ob alias egregias veritates ex his deducendas sese commendat.

\Article{340}
\textsc{Theorema.} \emph{Functio quaecunque aut est prima aut ex functionibus primis composita; posteriorique casu unico tantum modo e functionibus primis componi potest.}

\emph{Demonstr.} Nisi enim functio proposita $A$ sit prima, per aliam inferioris dimensionis $B$ dividetur. Si $B$ non est functio prima, per aliam $C$ inferioris gradus dividetur, itaque pergendo patet, tandem ad functionem primam deveniri, quoniam alias haec series foret infinita, quod, quoniam dimensiones perpetuo decrescunt, absurdum est. Iam si ultima functio prima sit $L$, haec omnes antecedentes metietur. Quare $A\equiv LA'$ eritque $A'$ inferioris dimensionis quam $A$. Quod iterum fiet $A'\equiv L'A''$ etc., patet, tandem ad functionem primam perveniri, adeoque $A$ erit $\equiv$ producto e functionibus primis $L,L',L''$ etc.\ $\pdot$Q.\,E.\,Pr.

Iam si etiam esset $A\equiv MM'M''$ etc. neque omnes $L,L',L''$ etc. eaedem cum omnibus $M,M',M''$ etc., eiiciantur eae, quae utrique seriei communes
\SourceLeaf{229}
sunt. Remaneantque $\lambda,\lambda',\lambda''\dotstwo$; $\mu,\mu',\mu'',\ldots$ eritque $\mu$ ad $\lambda,\lambda',\lambda''$ etc. prima, quare etiam ad productum $\lambda\lambda'\lambda''$ etc.; tamen esse debet
\[
\lambda\lambda'\lambda''\dotstwo\equiv\mu\mu'\mu''\dotstwo\text{ i. e. }\frac{\lambda\lambda'\lambda''\dotstwo}{\mu}\equiv\mu'\mu''\dotstwo\quad\text{Q. E. A.}
\]
